\section{Business Value For The Framework}

The empirical consciousness chapter changes the business status of the book in
two ways:
\begin{enumerate}[nosep]
  \item it gives the manuscript a real bridge to current biological research;
  \item it prevents the framework from sounding like an isolated productivity
        philosophy.
\end{enumerate}

It also improves later chapters. Once the reader accepts that state,
stability, and transition are serious scientific topics, planning, emergence,
and intervention design can be discussed as questions about regime management
rather than slogans about discipline.

\begin{discussion}
More concretely, the chapter upgrades the manuscript on four fronts. It
restores credibility to the state-language used throughout the book. It gives
later behavioural claims a cleaner boundary by stating what neuroscience does
and does not authorize. It improves the value of Chapter~\ref{ch:dynamics} by
showing that transition windows are not merely rhetorical. And it raises the
quality of the entire product by making the book look less like a private
framework and more like a disciplined teaching text with a real empirical
spine.
\end{discussion}
